\begin{python0}
from solutions import *; clear()
\end{python0}
\sectionthree{The [ ] operator}

Here's an example from the previous section.

\begin{consolethree}[escapeinside=||]
int a[3] = {4, 2, 42};
std::cout << a[0] << std::endl;

int x = a[1];

std::cout << x << std::endl;  
\end{consolethree}

Actually you can specify an index value using any \EMPHASIZE{integer expression}. Try this:

\begin{consolethree}[escapeinside=||]
int a[3] = {4, 2, 42};

int i = 0
std::cout << a[i] << std::endl;

int x = a[i + 1];
std::cout << x << std::endl;  
\end{consolethree}

In fact this is the reason why we have arrays ... Try this:

\begin{consolethree}[escapeinside=||]
int a[3] = {4, 2, 42};

for (int i = 0; i < 3; ++i)
{
    std::cout << a[i] << std::endl;
}
\end{consolethree}

The following code declares an array and populates it with integers entered by the user:

\begin{consolethree}[escapeinside=||]
int y[5] = {0};

for (int i = 0; i < 5; ++i)
{
    std::cin >> y[i];
} 
\end{consolethree}

Compare this to this code:

\begin{consolethree}[escapeinside=||]
int y0 = 0, y1 = 0, y2 = 0, y3 = 0, y4 = 0;

std::cin >> y0 >> y1 >> y2 >> y3 >> y4; 
\end{consolethree}

They more or less achieve the same thing, right? Now suppose instead of 5 variables, there are 1000 variables. The first version becomes:

\begin{consolethree}[escapeinside=||]
int y[1000] = {0};

for (int i = 0; i < 1000; ++i)
{
    std::cin >> y[i];
}
\end{consolethree}

Now let's see how much time you'd need to modify the \texttt{std::cin} statement if you have not heard of arrays:

\begin{consolethree}[escapeinside=||]
int y0 = 0, y1 = 0, y2 = 0, y3 = 0, y4 = 0, ...;

std::cin >> y0 >> y1 >> y2 >> y3 >> y4 >> y5 >> ...; 
\end{consolethree}

YIKES!!!

Get the point of arrays now?

In other words previously our for-loop allows us to iteratively execute
a block of statements. With the array, the for-loop allows us to execute
a block of statements on one variable from a huge list of variables one
at a time (if you think of an element of an array as a variable).
Previously our for-loop executes a bunch of statements on a single
variable (or maybe two or three). But here we apply the same computation
across a bunch of variables. POWERFUL IDEA!

In the next few sections, I'll show you how to apply
this principle in different scenarios.

One final thing before I show you different uses of the array.
It's always a good idea not to hardcode constants,
remember? The size of an array is not an exception. Therefore

\begin{consolethree}[escapeinside=||]
int y[1000] = {0};
for (int i = 0; i < 1000; ++i)
{
    std::cin >> y[i];
} 
\end{consolethree}

should be

\begin{consolethree}[escapeinside=||]const int SIZE = 1000;
int y[SIZE] = {0};

for (int i = 0; i < SIZE; ++i)
{
    std::cin >> y[i];
}
\end{consolethree}

For our example in this set of notes I will not use constants for array
sizes since the code is very short. But when the program is longer (such
as your assignments) you should use constants.

There's another way to avoid hardcoding the size of an
array. The function \texttt{sizeof()} will tell you the mount of memory
used by the computer to store the value of a variable. Try this:

\begin{consolethree}[escapeinside=||]int x = 0;
double y = 0.0;

char z = 'a';

int a[5] = {0};

std::cout << sizeof(x) << '\n'
          << sizeof(y) << '\n'
          << sizeof(z) << '\n'
          << sizeof(a) << '\n';

std::cout << sizeof(int) << '\n'
          << sizeof(double) << '\n'
          << sizeof(char) << '\n'; 
\end{consolethree}

So here's how you can determine the number of elements
in an array. Try this:

\begin{consolethree}[escapeinside=||]
int a[5] = {0};
double b[10] = {0};
char c[100];

std::cout << sizeof(a) / sizeof(int) << '\n';
std::cout << sizeof(b) / sizeof(double) << '\n';
std::cout << sizeof(c) / sizeof(char) << '\n'; 
\end{consolethree}

Get it?

One final thing. You should NEVER use an index value outside the bound
of your array. In other words if the size of your array a is 5, you MUST
not access \texttt{a[5]} or \texttt{a[1000]}; you should only access
\texttt{a[0]}, \texttt{a[1]}, \texttt{a[2]}, \texttt{a[3]}, and
\texttt{a[4].}

As mentioned earlier in the section on gotchas, you cannot assign
arrays:
\begin{consolethree}
int a[3] = {1, 2, 3};
int b[3];
b = a;               // BAD BAD BAD!!!
\end{consolethree}
This is what you have to do:
\begin{consolethree}
int a[3] = {1, 2, 3};
int b[3];
for (int i = 0; i < 3; ++i)
{
    b[i] = a[i];
}
\end{consolethree}
\begin{ex}
Declare an array of 10 integers; call this array \texttt{dice}. Using a for-loop, put random integers 1 to 6 into the array. Using another for-loop, print all the values in the array.
\end{ex}

